Los requisitos no funcionales describen las propiedades de calidad y restricciones que debe cumplir el sistema más allá de las funcionalidades que proporciona. Para su clasificación se ha tomado como referencia el modelo de calidad definido en el estándar ISO/IEC 25010 \cite{iso25010}.

La \autoref{tab:non-functional-requirements-executive} resume los requisitos no funcionales en formato ejecutivo, agrupados por característica de calidad para facilitar su lectura y trazabilidad.

{\scriptsize
\setlength{\tabcolsep}{3pt}
\renewcommand{\arraystretch}{1.14}
\newcommand{\RNFCategoryHeader}[1]{%
\hline\hline
\multicolumn{3}{|>{\RaggedRight\arraybackslash}p{0.91\linewidth}|}{\textbf{#1}}\\
\hline\hline
}
\begin{longtable}{|>{\RaggedRight\arraybackslash}p{0.18\linewidth}|>{\RaggedRight\arraybackslash}p{0.15\linewidth}|>{\RaggedRight\arraybackslash}p{0.58\linewidth}|}
\caption{Catálogo ejecutivo de requisitos no funcionales agrupado por características de calidad}
\label{tab:non-functional-requirements-executive}\\
\hline
\textbf{Característica} & \textbf{ID} & \textbf{Requisito no funcional sintetizado} \\
\hline
\endfirsthead
\hline
\textbf{Característica} & \textbf{ID} & \textbf{Requisito no funcional sintetizado} \\
\hline
\endhead
\RNFCategoryHeader{RNF1. Usabilidad}
Usabilidad & \texttt{RNF-U-01} & Proporcionar una interfaz intuitiva que permita realizar las tareas principales con un número reducido de interacciones. \\
\hline
Usabilidad & \texttt{RNF-U-02} & Adaptar la interfaz a móviles, tabletas y ordenadores, manteniendo una experiencia de usuario consistente. \\
\hline
Usabilidad & \texttt{RNF-U-03} & Presentar la información de forma clara y estructurada para facilitar la comprensión de datos operativos, comerciales y administrativos. \\
\hline
Usabilidad & \texttt{RNF-U-04} & Mantener criterios básicos de accesibilidad: etiquetas en formularios, nombres accesibles, contraste suficiente, texto legible y estructura revisable mediante WAVE. \\
\hline
\RNFCategoryHeader{RNF2. Eficiencia de rendimiento}
Rendimiento & \texttt{RNF-P-01} & Responder en tiempos razonables bajo uso normal, priorizando login, catálogo, clientes y operaciones sobre pedidos. \\
\hline
Rendimiento & \texttt{RNF-P-02} & Optimizar la carga de recursos para permitir un funcionamiento fluido también en dispositivos móviles con capacidad limitada. \\
\hline
\RNFCategoryHeader{RNF3. Compatibilidad}
Compatibilidad & \texttt{RNF-C-01} & Ser compatible con navegadores web modernos e informar mediante una ruta específica cuando el navegador no cumpla el soporte mínimo. \\
\hline
Compatibilidad & \texttt{RNF-C-02} & Ser utilizable desde los sistemas operativos móviles más habituales, especialmente Android e iOS mediante navegador. \\
\hline
Compatibilidad & \texttt{RNF-C-03} & Integrarse con servicios externos necesarios: almacenamiento de imágenes, geocodificación, rutas, QR, SMTP transaccional y Web Push. \\
\hline
\RNFCategoryHeader{RNF4. Fiabilidad}
Fiabilidad & \texttt{RNF-F-01} & Mantener disponibilidad suficiente para los usuarios y minimizar interrupciones del servicio durante la operación ordinaria. \\
\hline
Fiabilidad & \texttt{RNF-F-02} & Garantizar la conservación correcta de los datos almacenados, evitando pérdidas o inconsistencias de información. \\
\hline
\RNFCategoryHeader{RNF5. Seguridad}
Seguridad & \texttt{RNF-S-01} & Restringir el acceso a funcionalidades según el rol autorizado de cada usuario autenticado. \\
\hline
Seguridad & \texttt{RNF-S-02} & Proteger los datos mediante autenticación, autorización, trazabilidad de accesos y validación de operaciones sensibles. \\
\hline
Seguridad & \texttt{RNF-S-03} & Incorporar rate limiting y validación de sesión para reducir accesos abusivos o repetitivos sobre login y API. \\
\hline
\RNFCategoryHeader{RNF6. Mantenibilidad}
Mantenibilidad & \texttt{RNF-M-01} & Organizar el sistema en módulos y capas técnicas que faciliten su mantenimiento, evolución y ampliación. \\
\hline
Mantenibilidad & \texttt{RNF-M-02} & Permitir actualizar o modificar funcionalidades sin afectar innecesariamente al funcionamiento global de la aplicación. \\
\hline
\RNFCategoryHeader{RNF7. Portabilidad}
Portabilidad & \texttt{RNF-PO-01} & Ejecutar la aplicación en distintos dispositivos mediante navegador web, sin requerir instalaciones nativas específicas. \\
\hline
Portabilidad & \texttt{RNF-PO-02} & Permitir instalación como aplicación web mediante manifiesto e iconos, y habilitar notificaciones Web Push cuando el navegador lo soporte. \\
\hline
Portabilidad & \texttt{RNF-PO-03} & Ofrecer soporte mínimo ante fallos puntuales de conectividad mediante \textit{service worker} y caché del \textit{app shell}. \\
\hline
\end{longtable}
}
